\section{Assembly of the complete base}\label{sec:assembly}\label{assembly:section}

The local estimates are now available. We must realize infinitely many
analytic tests with one finite spatial dimension, preserve the earlier
tests when adding later wells, and obtain a complete global metric.
We choose the spatial data and join the local regimes. All tensor and
fixed-order derivative estimates in this section use operator or
multilinear norms in fixed affine coordinates. In particular, their
leading exponents do not include a trace over the spatial dimension.

\begin{proposition}\label{assembly:base}
For some fixed finite integer $n$, there is a smooth strictly
plurisubharmonic potential $\ell$ on $\mathbb C^n\times\mathbb C$
whose K\"ahler metric is complete and has nonpositive real sectional
curvature. The potential satisfies every hypothesis of
Proposition~\ref{grid:criterion}, for a countable dense set of centers
in the unit disk and every spatial Taylor order. Thus all bounded
holomorphic functions on its half-plane tube are constant.
\end{proposition}

At each spatial point the fiber potential is recovered from
\[
 F(\log\tau,s)=\log|y-\zeta(s)|^2+\psi(s)
\]
with the left normalization of Section~\ref{sec:fiber}. The radial
function $B$ supplies the spatial metric. The twist $\psi$ starts with
a radial background $\Psi$: polynomial wells make selected fibers wide,
and additional radial plurisubharmonic terms control the derivatives
of the severity. A term
$\operatorname{Re}(c_{\rm sh}(s)(y-\zeta(s)))$ prescribes the sheared
disk directions. We choose these data in increasing spatial radius.
The polynomial wells are the only later additions that extend into
earlier regions, so their inward tails must be reserved in advance.

The proof occupies this section. A cycle starts on a capped spatial
interval, raises the severity, performs one test, and raises severity
again on a post-test ramp. At constant severity it then accelerates the
spatial metric, enters the aligned regimes, changes center, and restores
the spatial parameters. Finally it increases severity while the
normalized spatial geometry is fixed, until the severity can pay for
the next capped interval.
\begin{figure}[tbp]
\centering
\begin{tikzpicture}[
  stage/.style={draw=black!55,rounded corners=2pt,fill=black!3,
    text width=3.65cm,minimum height=1.27cm,align=center,font=\small},
  flow/.style={-{Latex[length=2mm]},thick,draw=black!65},
  note/.style={font=\footnotesize,align=center}]
\node[stage] (cap) at (0,0) {Cap and coast\\
  $\lambda=1$, $H\le\sigma$};
\node[stage] (trigger) at (4.6,0) {Severity trigger\\
  raise $E$ to $E_*$};
\node[stage] (test) at (9.2,0) {Polynomial test\\
  fixed center and severity};
\node[stage] (post) at (9.2,-2.05) {Post-test ramp\\
  reach $E_{\rm post}$};
\node[stage] (accelerate) at (4.6,-2.05) {Accelerate growth\\
  enter aligned regime};
\node[stage] (grow) at (0,-2.05) {Grow $H$, then $\lambda$\\
  $E=E_{\rm post}$};
\node[stage] (center) at (0,-4.1) {Widen, move, restore\\
  hold $\kappa=p=1$};
\node[stage] (reset) at (4.6,-4.1) {Return $\lambda$ to $1$\\
  then lower $H$ to $\sigma$};
\node[stage] (catchup) at (9.2,-4.1) {Severity catch-up\\
  before the next cap};
\draw[flow] (cap)--(trigger);
\draw[flow] (trigger)--(test);
\draw[flow] (test)--(post);
\draw[flow] (post)--(accelerate);
\draw[flow] (accelerate)--(grow);
\draw[flow] (grow)--(center);
\draw[flow] (center)--(reset);
\draw[flow] (reset)--(catchup);
\draw[flow] (catchup.south)--++(0,-.48)
  node[below,note]{Next stage: cap and coast};
\end{tikzpicture}
\caption{One stage, in increasing spatial radius. The test is installed
while the spatial growth is capped. The intervening growth and center
change occupy finite clock intervals. After the reset, severity must
catch up with the accumulated spatial scale before the next capped
regime can begin. The arrows indicate order, not elapsed radius.}
\label{fig:stage-clock}
\end{figure}


We retain the radial notation of Lemma~\ref{radial:odes}:
\[
 \begin{gathered}
 t=|s|^2,\quad \nu=\log t,\quad a=B'(t),\quad u_s=\log a,\\
 Y=\frac{du_s}{d\nu},\quad a_\psi=\Psi'(t)=a/\lambda,\quad
 H=\frac{Y}{ta_\psi},\quad p=\frac{d\log a_\psi}{du_s}.
 \end{gathered}
\]
Here $H$ is a scalar; the polynomial system below is denoted by
$\mathcal H$. The evolution equations are
\begin{align}
 \frac{dY}{du_s}&=\kappa(1+Y),&
 \frac{d\nu}{du_s}&=Y^{-1},&
 \frac{d\log\lambda}{du_s}&=1-p,\label{assembly:clocks}\\
 \frac{d\Psi}{du_s}&=H^{-1},&
 \frac{d\log H}{du_s}
 &=\kappa(1+Y^{-1})-p-Y^{-1}.
 \label{assembly:H-clock}
\end{align}
The admissibility and derivative conditions on $\kappa,p$ in the radial
lemmas will always be imposed.

\subsection{Wells and the order of constants}

Fix a sequence $(\zeta_j,h_j)$ in which every pair from a fixed countable
dense subset of $\{|y|<1\}$ and $\mathbb Z_{\geq0}$ occurs infinitely
often. The parameters of a test will be
\begin{equation}\label{assembly:test-parameters}
 E_*=B_*c,\qquad \eta=e^{-c},\qquad
 N=\lfloor e^{\gamma E_*}\rfloor,\qquad
 l=\frac N2\log(t/R^2),\qquad S=Re^{c/N}.
\end{equation}
The constants $B_*,\gamma$ are fixed below. At each stage $R$ is fixed
before $c$. For a degree-$N$ system $(\mathcal H,p_N)$ from
Lemma~\ref{grid:polynomials}, define
\begin{equation}\label{assembly:well}
 W_j(s)=d\log
 \frac{|\mathcal H(s/R)|^2+|p_N(s/R)-1|^2+e^{-2c}}
 {1+e^{-2c}}.
\end{equation}
The polynomial $p_N$ is unrelated to the radial control $p$.
The logarithmic sum-of-squares formula shows that $W_j$ is smooth
and plurisubharmonic.

\begin{lemma}\label{assembly:well-bounds}
There is an absolute exponent $C_{\rm pol}$ such that, for $t\geq1$
and $1\leq r\leq4$,
\begin{equation}\label{assembly:well-raw}
 |\partial^{[r]}W_j|
 \leq C(n,d)\exp\{C_{\rm pol}(\max(l,0)+c+\log N)\}.
\end{equation}
There is a constant $L_n<\infty$ such that, after increasing the degree
threshold,
\begin{equation}\label{assembly:well-tail}
 |\partial^{[r]}W_j|\leq N^{L_n}
 \qquad(1\leq r\leq4,\ l\geq L_n\log N).
\end{equation}
Both statements hold throughout the indicated outer ranges.
\end{lemma}

\begin{proof}
Homogeneity and the spherical upper bound give
$|(\mathcal H,p_N)(s/R)|\leq Ne^l$. Cauchy estimates at scale
$c_0\sqrt t/N$ bound its derivatives of order at most four by this
quantity times a fixed power of $N$, since $t\geq1$. In a differentiated
logarithm there are only a fixed number of factors and denominator
powers; the denominator in \eqref{assembly:well} is at least $e^{-2c}$.
This proves \eqref{assembly:well-raw} with an absolute exponent;
fixed-dimensional constants and $d$ affect only its prefactor.

The spherical lower bound also gives
$|(\mathcal H,p_N)(s/R)|\geq e^lN^{-C_n}$. Once $l$ exceeds a sufficiently
large multiple of $\log N$, subtraction of $(0,1)$ leaves norm at least
$e^lN^{-C_n}/2$. The factors $e^l$ in derivatives of the logarithm can
then be canceled. The remaining costs are fixed powers of $N$, with
exponents depending on the already fixed $n$. Increase $L_n$ to pay
them and the degree threshold to pay the prefactor.
\end{proof}

For each stage we also introduce smooth nonnegative radial
plurisubharmonic functions $Z_{j,\rm in}$ and $Z_{j,\rm out}$. They
vanish before their stage and eventually become affine nondecreasing
functions of $\nu$. The center changes add offsets that vanish before
their change interval and are constant afterwards. We seek the final
twist in the form
\begin{equation}\label{assembly:twist}
 \psi=\Psi+\sum_j W_j+\sum_j(Z_{j,\rm in}+Z_{j,\rm out})
       +\text{center-change offsets}.
\end{equation}
The center is constant except on its change intervals. There is one
fiber profile at each spatial point: it is $q_0(x,E(s))$ except during
the coordinated switches of Proposition~\ref{center:change}.
Shears occur only in the disjoint test collars. The recursion below
will make the well sum converge smoothly on compact spatial sets;
the other added terms are locally finite.

The constants must provide the three outputs of
Proposition~\ref{grid:criterion}: a fixed-height spatial ball, small
unsheared disks on its boundary, and large prescribed sheared disks at
the nodes. Their exponents determine $k$ and hence the dimension.
Accordingly we first choose all geometric constants independently of
$n$, and only then use stagewise size choices to pay dimensional
prefactors.

We give the full order of choices. First choose the profile constants,
including $d_0,K_0$ for the center change, then $\epsilon,\mu$ for
Proposition~\ref{quiet:profile}, the shear exponent in
Proposition~\ref{quiet:shear}, and the tolerance $\delta$ for
Proposition~\ref{radial:aligned}. These choices are independent of $n$.
Use control transitions whose first two derivatives in $\log Y$ have
absolute bounds. Let $C_{\rm rad}$ dominate the absolute exponents
needed for radial raw jets through order four. Choose
\begin{equation}\label{assembly:first-constants}
 \epsilon_1\ll\epsilon/(1+C_{\rm rad}),\qquad
 C_\uparrow\gg(1+C_{\rm rad})/\epsilon.
\end{equation}
Choose next
\begin{equation}\label{assembly:second-constants}
 m_z\geq10(C_{\rm pol}+1),\qquad
 B_*\geq\frac{C_*(1+m_z+C_{\rm pol}+C_{\rm rad})}{\epsilon_1},
\end{equation}
where the absolute $C_*$ is large enough for the finitely many raw
exponent comparisons below. Choose $C_s$ larger than the quiet shear
severity exponent, and $C_w$ large enough to pay the shell loss and
four derivatives of its cutoffs; $C_w>5k_*+20$ suffices for these latter
costs after increasing the other fixed constants if necessary. Here
$k_*$ is the absolute separation exponent of the polynomial grids.
Choose
\begin{equation}\label{assembly:gamma-kb}
 0<\gamma\ll\epsilon_1,\qquad
 \gamma B_*\leq\min\{.005,(8C_w)^{-1}\},\qquad
 k_b>\max\{8/\gamma,100/\epsilon\}.
\end{equation}
Choose $M_b$ sufficiently large for the ratios in
Lemma~\ref{assembly:quiet-check}, and then $\sigma>0$ sufficiently small
for that lemma and Lemma~\ref{assembly:reset}. These are choices made
from fixed constants; in particular
$\sigma M_bk_b^2\ll1$, and $\sigma$ pays all quiet smallness requirements
involving $C_\uparrow$. Severity thresholds can depend on $\sigma$.
Finally take
\begin{equation}\label{assembly:depth}
 d\geq8(M_bk_b+C_s+1)B_*.
\end{equation}
All constants so far are independent of dimension.

Set
\begin{equation}\label{assembly:test-exponents}
 m_0=d/2,\qquad p_0=2d,\qquad
 p_1=C_sB_*+C_wd\gamma B_*.
\end{equation}
Then $p_1\leq d/4<m_0$. Choose an integer $k$ with
$.99k>p_0+p_1+1$, and only now fix a finite $n$ sufficiently large that
\begin{equation}\label{assembly:dimension}
 \sum_{r=1}^{k-1}\frac{r^{n-1}}{(n-1)!}<.01.
\end{equation}
Such a finite $n$ exists because $k$ is fixed. The constants $C_n,L_n$
and all dimensional prefactors will be paid by the stagewise choice
of large $c$; they do not change \eqref{assembly:test-exponents}.

\subsection{Recursive convention and capped intervals}

Every estimate for an already configured region is made with an
additional allowance of one for the absolute value and spatial
jets through order four of all future wells. Let $\nu_j^{\rm sat}$
denote the designated endpoint of severity saturation at stage $j$,
and let
\[
 K_j=\{s:|s|^2\leq e^{\nu_j^{\rm sat}}\}
\]
be the entire inner spatial ball ending there. The testing annulus lies
beyond this ball. At stage $j$ we impose
\begin{equation}\label{assembly:diagonal}
 \sup_{s\in K_j}|\partial^{[r]}W_j(s)|\leq2^{-j},
 \qquad 0\leq r\leq j+4.
\end{equation}
Each new trigger begins beyond every completed stage, so every already
configured region lies in every later $K_j$. The omitted well sum
therefore has absolute value and spatial jets through order four
bounded by $\sum_j2^{-j}\leq1$ there. Its Hessian is also positive
semidefinite. Feasibility is proved below. This convention reserves the
entire future allowance before selecting any future coefficients. We apply
the local all-fiber lemmas to the final twist with these allowances;
we do not infer a uniform curvature statement from small perturbations
on compact subsets of the total space.

Initially put
\begin{equation}\label{assembly:initial-data}
 a=a_\psi=a_{\rm init}e^t,\quad Y=t,\quad
 \lambda=p=1,\quad\kappa=(1+Y)^{-1},
\end{equation}
with $a_{\rm init}$ so large that $H\leq\sigma$, and choose $B,\Psi$
with equal initial constants. Up to a positive scale and an additive
constant, this initial spatial potential is the exponential radial
example of \citet[Section~3]{Seshadri2006}. Take center $\zeta_1$ and a constant
severity sufficiently large for the initial compact region and the
tail allowance. The radial curvature formula is strictly positive
on nonzero rank-two Hermitian tests at $t=0$, by its limit or by
direct differentiation of $a_{\rm init}e^t$. Thus the quiet criterion
also applies at the spatial origin.

Before every trigger the radial data will have the form
\begin{equation}\label{assembly:cap}
 \lambda=p=1,\qquad
 \kappa=\frac{C_\kappa}{1+Y},\qquad
 1\leq C_\kappa\leq1+Y,\qquad H\leq\sigma,
\end{equation}
with $C_\kappa$ constant on that interval. Then
$dY/d\nu=C_\kappa Y$, so the cap can be continued through any
prescribed finite interval of $\nu$ without a pole. All previously
generated nonradial twist terms have bounded positive-order Euclidean
jets on the entire outer range: use \eqref{assembly:well-tail} and
the affine $Z$ tails. Hence a cap can also be continued until
$a=e^{u_s}$ dominates any prescribed bound for those old jets.

Choose the next trigger start $\nu_B$ with
$Y_B=Y(\nu_B)\gg C_\kappa$, $e^{\nu_B}\geq1$, and severity constant
there at a value satisfying
\begin{equation}\label{assembly:trigger-initial}
 E_0\geq C_\uparrow u_s(\nu_B).
\end{equation}
The reset below provides this inequality and any additional fixed
threshold. Put
\begin{equation}\label{assembly:trigger-radius}
 \xi=Y_B(\nu-\nu_B),\qquad
 \nu_*=\nu_B+Y_B^{-1},\qquad R=e^{\nu_*/2}.
\end{equation}
These data precede the choice of $c$. Retain the capped radial data
through
\begin{equation}\label{assembly:post-length}
 l=T_*+4,\qquad T_*=10\{c+(1+L_n)\log N\}.
\end{equation}
For large $c$ this lies inside $\nu\leq\nu_B+2/Y_B$, since
$(T_*+4)/N\to0$. Throughout this fixed continuation
$Y\asymp Y_B$ and $u_s=u_s(\nu_B)+O(1)$, with constants allowed
to depend on the already fixed cap data.

\begin{lemma}\label{assembly:height}
Before choosing $c$ at stage $j$, one can fix $D_j\geq j$ such that,
for every sufficiently large permitted $c$ and every final choice of
the reserved tails,
\begin{equation}\label{assembly:height-bound}
 \ell(s,\zeta_j)\leq D_j-2\qquad(|s|\leq Re^{c/N}).
\end{equation}
\end{lemma}

\begin{proof}
On the spatial region following the previous center-change interval,
the center is $\zeta_j$ and $\ell(s,\zeta_j)=B(s)$, including on shear
supports. The same identity holds on the initial cap at the first
stage. Thus the bound there follows from the fixed cap continuation
up to $\nu_B+2/Y_B$.
The earlier spatial region is a fixed compact set $K$ with already
chosen profiles and centers; future stages do not change $q(x,s)$
on $K$. The reservation gives $|\psi-\psi_{\rm old}|\leq1$ there,
including the current and all later wells. Both centers are in the unit disk, so
$|\zeta_j-\zeta(s)|\leq2$. Choose
$A>\log4+\sup_K\psi_{\rm old}+1$.
For each $s\in K$, $F(x,s)\to\infty$ as $x\to\infty$.
Continuity, a finite cover, and monotonicity in $x$ give a common
finite $X$ with $F(X,s)>A$ on $K$. Inverting the fiber relation
therefore gives $\tau\leq e^X$. Since $q\geq1$, the radial potential
is bounded by $\sup_K B+e^X$; the fixed shear adds at most
$2\sup_K|c_{\rm sh}|$. At center coincidences this argument is
understood by the smooth limit $\tau=0$. These bounds are independent
of $c$ and of the future wells. Add two and increase to at least $j$.
\end{proof}

We always require
\begin{equation}\label{assembly:height-rate}
 c\geq j(D_j+1)
\end{equation}
in addition to the finitely many current thresholds. Thus $D_j\to\infty$
and $D_j/c_j\to0$.

\subsection{The severity trigger}

Choose a smooth convex nondecreasing $h_b$ on $(-\infty,1)$, zero
for $\xi\leq0$, equal to $-\log(1-\xi)$ up to an affine function near
$1$. To construct it, cut on the nonnegative second derivative
$(1-\xi)^{-2}$ smoothly and integrate twice from zero.
Use
\begin{equation}\label{assembly:trigger}
 E=E_0+k_bh_b(\xi),\qquad Z_{j,\rm in}=M_bk_b^2h_b(\xi)
\end{equation}
until severity is near $E_*$. Smoothly saturate $E$ at $E_*$ over
a unit interval of its argument, so that $E=E_*$ by
$k_bh_b=E_*-E_0+1$. Derivatives of the saturation function in that
argument are bounded. Keep the unsaturated convex $Z_{j,\rm in}$
slightly longer, and once $k_bh_b\geq E_*-E_0+2$, cut off its
second derivative nonnegatively over half the remaining gap to
$\xi=1$. Continue it affinely in $\nu$ thereafter. This is a smooth
convex nondecreasing function. Such a function of $\log|s|^2$,
zero near the origin, is a smooth radial plurisubharmonic function:
its transverse and radial Hessian entries are its first and second
$\nu$ derivatives divided by $t$.

Near this cutoff the gap $1-\xi$ is comparable, with fixed stagewise
factors, to $e^{-(E_*-E_0)/k_b}$. The final slope is bounded by
\begin{equation}\label{assembly:trigger-slope}
 C M_bk_b^2Y_Be^{(E_*-E_0)/k_b}.
\end{equation}
The value there is $M_bk_b(E_*-E_0+O(1+k_b))$. Multiplication
of \eqref{assembly:trigger-slope} by the remaining gap to $\nu_*$
costs only a fixed $O(M_bk_b^2)$, and its product with $4c/N$
tends to zero. Thus for large $c$,
\begin{equation}\label{assembly:trigger-value}
 Z_{j,\rm in}\leq2M_bk_bE_*\qquad(|l|\leq2c).
\end{equation}
Through saturation, its derivatives and those of $E$ through order
four have bounds given by fixed prefactors times
$(1+Y_B)^4 e^{4(E-E_0)/k_b}$.
The bound with $E_*$ also holds for positive-order derivatives of
the affine tail.

The well satisfies \eqref{assembly:diagonal} after increasing $c$.
Indeed throughout severity saturation,
\begin{equation}\label{assembly:activation-gap}
 l\leq-C(Y_B,E_0)N e^{-E_*/k_b}.
\end{equation}
The absolute value on the right grows exponentially in $E_*$,
because $\log N\sim\gamma E_*$ and $\gamma>4/k_b$.
Homogeneity makes the polynomial evaluations on the whole inner
ball exponentially small in this quantity. Cauchy on a collar
of scale $R/N$ gives the same conclusion for any fixed derivative
order, with only powers of $N$ as further costs. The logarithm in
\eqref{assembly:well} is then close to zero with all these
derivatives. We can therefore enforce the bound $2^{-j}$ for its
value and derivatives through order $j+4$. This decay pays all
fixed exponential-in-$c$ costs needed at stage $j$.

\begin{lemma}\label{assembly:quiet-check}
The trigger satisfies the hypotheses of
Proposition~\ref{quiet:profile}, uniformly for every $\tau\geq0$.
After saturation the same criterion applies through $l=4c$, including
the initial part of the post-test term defined below.
\end{lemma}

\begin{proof}
For raw estimates use Euclidean axes. Here $t\geq1$, $Y\asymp Y_B$,
and $\log Y_B\leq C u_s(\nu_B)$; the last estimate follows by
integrating $d\log Y/du_s=\kappa(1+Y^{-1})\leq2.2$ for $Y\geq1$.
Differentiation of $a'=aY/t$ bounds the required radial metric jets
by $e^{C_{\rm rad}u_s}$. Their inverses cost no more. The radial
curvature estimate gives
\begin{equation}\label{assembly:euclidean-curvature}
 C^V(T,T)\geq c\kappa H|T|_V^2
 \geq c(a/t)|T|^2,
\end{equation}
since $V\geq aI$ and
$\kappa H=C_\kappa Y/((1+Y)ta)\geq c/(ta)$.
The initial severity pays these bounds, the old data and the future
allowance. The trigger cost $e^{4E/k_b}$ fits the quiet raw exponent.
The current well is small before saturation. Also $G\geq V$, since
all nonconstant extra twist terms in this interval are
plurisubharmonic.

For the covariant condition, write
\[
 \mathcal B(\nu)=B(e^\nu),\qquad
 \mathcal U(\nu)=\Psi(e^\nu)+Z_{j,\rm in}(e^\nu),
\]
and use primes for $\nu$ derivatives within this proof. On the
severity-varying interval,
\begin{align}
 \frac{\mathcal U''}{\mathcal U'},\
 \frac{|\mathcal U'''|}{\mathcal U''},\
 \frac{\mathcal B''}{\mathcal B'},\
 \frac{|\mathcal B'''|}{\mathcal B''}
 &\leq C\sqrt{\mathcal U''},
 \label{assembly:connection-ratios}\\
 \frac{|E'|}{\mathcal U'},\
 \frac{|E'|}{\sqrt{\mathcal U''}},\
 \frac{|E''|}{\mathcal U''},\
 \frac{|E'''|}{(\mathcal U'')^{3/2}}
 &\leq\varepsilon_b,
 \label{assembly:small-ratios}
\end{align}
where $\varepsilon_b$ can be made as small as required.
Away from the pole, $h_b$ and its relevant derivatives are bounded,
whereas
$\mathcal B'=Y/H$ and $\mathcal B''=Y(1+Y)/H$.
The choice $\sigma M_bk_b^2\ll1$ makes these background terms
dominate. Near the pole let $D_b=Y_B/(1-\xi)$. Then
\[
 Z_{j,\rm in}'\asymp M_bk_b^2D_b,\quad
 Z_{j,\rm in}''\asymp M_bk_b^2D_b^2,\quad
 |Z_{j,\rm in}'''|\leq C M_bk_b^2D_b^3,
\]
and $|E^{(r)}|\leq Ck_b^rD_b^r$ for $1\leq r\leq3$, including
saturation. These prove \eqref{assembly:connection-ratios} and
make \eqref{assembly:small-ratios} small by first enlarging $M_b$
and then decreasing $\sigma$. The fixed intervening interval of
$\xi$ is covered by the first estimates. We also have
$\mathcal U''\geq1$.

At fixed $\tau$, the radial part of the Schur metric is the Hessian of
\[
 J_h=\mathcal B+\tau(\mathcal U-\bar F(x,E)),
\]
where the bar denotes averaging from zero to $\tau$ in the $\tau$
variable. Bounded severity derivatives of $\bar F$ and
\eqref{assembly:small-ratios} show that the first two derivatives
of $\mathcal U-\bar F$ are positively comparable with those of
$\mathcal U$ and its third derivative is bounded by
$C(\mathcal U'')^{3/2}$.
Use affine axes normalized by the Hessian of $\mathcal U$, with
radial radius $\sqrt{t/\mathcal U''}$ and transverse radii
$\sqrt{t/\mathcal U'}$.
At a radial point, the only nonzero holomorphic connection types
for a radial Hessian metric are radial--radial to radial and
radial--transverse to transverse. Differentiation gives the bounds
\begin{equation}\label{assembly:connection-exact}
 \frac{|J_h'''/J_h''-1|}{\sqrt{\mathcal U''}},
 \qquad
 \frac{|J_h''/J_h'-1|}{\sqrt{\mathcal U''}}.
\end{equation}
They are bounded by \eqref{assembly:connection-ratios}. No estimate
of a radial/transverse condition number is used.

Write the remaining Hessian as $\tau G_e$, with $G_e\geq0$.
Its size and first jets in these axes are bounded: old Euclidean
bounds are dominated by the large initial $a$, and the current
and future wells have the reserved absolute bounds. The normalizing
radii are at most $a^{-1/2}$ because this background is at least
$V\geq aI$. If $h_r$ denotes the radial Schur metric, the exact
connection difference is, with the usual index placement,
\begin{equation}\label{assembly:connection-perturbation}
 \Gamma^h-\Gamma^{h_r}
 =\tau(\partial G_e-\Gamma^{h_r}G_e)(h_r+\tau G_e)^{-1}.
\end{equation}
Since $h_r\geq c\tau I$, $\tau h^{-1}$ is uniformly bounded for
all $\tau>0$. Thus the full connection is uniformly bounded as well.

In these axes the first derivative of $E$ has norm
$|E'|/\sqrt{\mathcal U''}$; the complex Hessian entries are bounded
by $|E'|/\mathcal U'$ and $|E''|/\mathcal U''$, and the only pure
second derivative has size $|E''-E'|/\mathcal U''$.
Consequently subtracting the connection times the first derivative
gives a small $\nabla_h^{2,0}E$. Since the actual $G$ is at least the
normalizing background, the same smallness holds in $G$ units.
This proves all quiet smallness conditions uniformly in $\tau$.

After saturation, severity derivatives vanish. Through $l=4c$ the
well's raw costs are bounded by $e^{C(c+\log N)}$; the post term
below has cost at most a fixed multiple of $N^4e^{2m_zc}$.
Both fit $\epsilon E_*$ by \eqref{assembly:second-constants}, while
the inverse and radial curvature estimates persist.
\end{proof}

\subsection{Testing disks and the post ramp}

Choose $c$ sufficiently large that the degree in
\eqref{assembly:test-parameters} meets the grid accuracy required for
$h_j$, and $N>h_j$. At every good lift $v$, use a smooth cutoff
supported in $|s-Rv|<2RN^{-k_*}$ and equal to one on
$|s-Rv|\leq RN^{-k_*}$. Assign on the inner ball
\begin{equation}\label{assembly:shear-amplitude}
 c_{\rm sh}=\alpha b_v,\qquad
 \alpha=\exp(-C_sE_*-C_wd\log N),
\end{equation}
and multiply by the cutoff in the collar. The phases have modulus
one, are independently assignable by projective direction, and
repeat on every lift of that direction. The outer balls are disjoint.
A cutoff $\chi(|s-Rv|^2/(RN^{-k_*})^2)$ has fixed-order multilinear
derivative bound $C_r(RN^{-k_*})^{-r}$ with absolute $C_r$.
At most one contributes at any point, so there is no factor for
the number of nodes or dimension.

On the derivative supports, the shell estimate gives
\begin{equation}\label{assembly:shear-lower}
 \psi\geq-2d(k_*+11)\log N-C_{\rm old}.
\end{equation}
Old values are bounded on the fixed cap continuation, future values
have the allowance, and the trigger term is nonnegative.
These balls have $|l|\leq\log N+1$ and lie beyond severity
saturation. First and pure second derivatives of $\psi$ also fit
the quiet raw bounds; the extra background first-derivative cost
is at most $a\sqrt t$. The constants $C_s,C_w$ pay respectively
the severity exponent of Proposition~\ref{quiet:shear} and half the
loss \eqref{assembly:shear-lower} together with the four cutoff
derivatives. Large $c$ pays fixed old factors. Thus that proposition
applies for all fiber radii with metric comparison loss at most
$1/2$. Where a shear coefficient is constant, its potential term
is pluriharmonic.

At a node the new well is $-2dc+O(1)$; on $|s|=S$ it is at most
$2d(c+\log N)+O_d(1)$. All other values on the fixed continuation
are bounded independently of $c$, except the trigger term.
The post term is still zero. Hence
\begin{align}
 \psi_{\rm node}&\leq-2dc+2M_bk_bB_*c+C_{\rm old},
 \label{assembly:node-twist}\\
 \psi_{\rm sphere}&\leq2d(c+\log N)+2M_bk_bB_*c+C_{\rm old,d}.
 \label{assembly:sphere-twist}
\end{align}
For a small fixed profile-dependent $c_0>0$, the disk
\begin{equation}\label{assembly:fiber-disk}
 |y-\zeta_j|<c_0e^{-(\psi+E_*)/2}
\end{equation}
has $\tau\leq e^{-E_*}$ and radial potential increment at most one.
Indeed, by the axis estimates in Lemma~\ref{quiet:scalar}, at
$\tau=e^{-E_*}$ the normalized relation gives $F=-E_*+O(1)$ uniformly: integrate $1/q-1=O(e^{E_*}\tau)$
up to that value of $\tau$. Monotonicity proves the disk assertion.

At a node the holomorphic shear
$w\mapsto w+i\alpha b_v(y-\zeta_j)$ has imaginary increment
$\operatorname{Re}(\alpha b_v(y-\zeta_j))$, exactly the added
potential. Writing $z=y-\zeta_j$, the clearance on the disk
\eqref{assembly:fiber-disk}, over any $w$ with
$\operatorname{Im}w=D_j$, is therefore
\[
 \begin{split}
 &\operatorname{Im}(w+i\alpha b_vz)
       -\ell(Rv,\zeta_j+z)\\
 &\qquad=D_j-\ell_{\rm rad}(Rv,\zeta_j+z)
       \geq D_j-B(Rv)-1\geq1.
 \end{split}
\]
The last inequality uses \eqref{assembly:height-bound} on the center
axis. Thus the shear consumes none of the height margin, and the
sheared disk lies in the tube. Its logarithmic radius is at least
\[
 (d-M_bk_bB_*-B_*/2)c-C_{\rm old}
 \geq(13d/16)c-C_{\rm old},
\]
which exceeds $m_0c$ for large $c$.
On the sphere there is no shear, and Cauchy's inequality on
\eqref{assembly:fiber-disk} bounds the transverse derivative
of a function of modulus at most one by
\[
 \exp\{[d(1+\gamma B_*)+M_bk_bB_*+B_*/2]c+C_{\rm old,d}\}
 \leq e^{p_0c}.
\]
Also $\alpha^{-1}\leq e^{p_1c}$.
These are precisely the constants \eqref{assembly:test-exponents},
with strict margins to absorb the prefactors.
Choose the phases supplied by Proposition~\ref{grid:criterion}.
All geometric estimates hold for arbitrary phases.

For the post ramp, add a smooth convex nondecreasing function
$Z_{j,\rm out}$ of $l$, zero through $2c$, and equal to
$e^{m_z(l-2c)}$ up to affine terms from $2c+2$ onward.
This is obtained by cutting on its positive second derivative and
integrating twice. Continue through $T_*+2$, then cut off that
second derivative nonnegatively so that the function is affine
after $T_*+3$. Independently turn on the severity slope $B_*$
in $l$ over $[4c,4c+1]$, and turn it off over
$[T_*,T_*+1]$. Write $E_{\rm post}$ for the final value.
For large $c$,
\begin{equation}\label{assembly:post-severity}
 E\geq B_*(c+l)/10\quad(2c\leq l\leq T_*+4),\qquad
 E_{\rm post}\geq B_*T_*/2.
\end{equation}

\begin{lemma}\label{assembly:post}
The quiet criterion holds throughout the post ramp.
Beyond $l=T_*+4$, the sum of all already generated additive twist
terms other than $\Psi$ and the constant offsets, and each such
term individually, has Euclidean positive-order jets through order
four bounded on the whole outer range by
\begin{equation}\label{assembly:outer-bound}
 D_{\rm out}=\exp(\epsilon_1E_{\rm post}/3).
\end{equation}
\end{lemma}

\begin{proof}
Each derivative of a function of $l$ costs a fixed power of $N$.
Thus $Z_{j,\rm out}$ has positive-order bounds through order four
of size $C N^4e^{m_z\max(0,l-2c)}$. These and
\eqref{assembly:well-raw} fit the raw quiet exponent by
\eqref{assembly:second-constants} and \eqref{assembly:post-severity}.
Other radial and old-data bounds remain fixed on the cap
continuation. Through $l=4c$, severity is constant and the preceding
lemma applies.

Where severity varies, take $\mathcal U=\Psi+Z_{j,\rm out}$.
Before division by $t$, its transverse and radial Hessian entries
have sizes
\[
 m_zN e^{m_z(l-2c)},\qquad
 (m_zN)^2e^{m_z(l-2c)}.
\]
For $l\geq4c$, $m_z(l-2c)$ exceeds
$C_{\rm pol}(l+c+\log N)$ by a fixed positive multiple of $c+l$.
Thus this Hessian dominates the current well and its first
Hessian jets in normalized axes.
The affine $Z_{j,\rm in}$ tail has coefficient at most a fixed
factor times $e^{E_*/k_b}\ll N$; old terms and future tails are
covered as before. Moreover
$\mathcal U''/\mathcal U',|\mathcal U'''|/\mathcal U''\leq Cm_zN$,
while $\sqrt{\mathcal U''}\gtrsim
m_zN e^{m_z(l-2c)/2}$.
Severity derivatives of orders $r\leq3$ are at most $CB_*N^r$.
Thus the connection ratios hold and all severity ratios are small.
Equations~\eqref{assembly:connection-exact} and
\eqref{assembly:connection-perturbation} again prove the pure
Hessian condition uniformly for all $\tau$. Every extra Hessian
is positive semidefinite in this interval.

Past $T_*+4$ the new well satisfies \eqref{assembly:well-tail},
and the $Z$ terms are affine in $\nu$. Their positive-order
Euclidean jets are bounded on the entire outer range by their
endpoint bounds up to fixed constants. Their logarithmic costs
are at most
$C(1+m_z)T_*+4\log N+C_{\rm old}$ with an absolute $C$.
The absolute constant $C_*$ in \eqref{assembly:second-constants}
was chosen to pay this inequality too.
Equation~\eqref{assembly:post-severity} then implies
\eqref{assembly:outer-bound}, after increasing $c$ for the finite
sum of old bounds, the number of generated terms, and dimensional
prefactors. The dimension-dependent $L_n$
multiplies $T_*$ on both sides and changes no leading comparison.
\end{proof}

\subsection{Acceleration, center change, and reset}

At constant $E_{\rm post}$, increase $\kappa$ from the cap value
to $\kappa_*=1.08$ on a fixed-factor interval of $Y$ after
$l=T_*+4$, keeping $\lambda=p=1$. The transition can preserve
$Y\kappa_Y\geq-\kappa$: multiply the difference between the cap
value and $\kappa_*$ by a nondecreasing cutoff in $\log Y$.
The upward derivative only improves that inequality.
The $u_s$ increment is bounded in terms of the fixed prior $Y_B$,
independently of large $c$. Keep $\kappa=\kappa_*$ until
\begin{equation}\label{assembly:aligned-entry}
 u_s=\epsilon_1E_{\rm post}.
\end{equation}
For large $c$ this endpoint lies after the transition and has
$H\geq\sigma$. Indeed, with $p=1$,
$d\log H/du_s=(\kappa_*-1)(1+Y^{-1})\geq.08$,
whereas the initial values at the end of the cap have fixed
stagewise bounds as $c\to\infty$.

Until \eqref{assembly:aligned-entry}, apply the quiet criterion.
The radial raw jets cost $e^{C_{\rm rad}u_s}$, the extra twist
jets are bounded by \eqref{assembly:outer-bound}, and the curvature
lower bound follows as in \eqref{assembly:euclidean-curvature}.
Severity derivatives vanish. The choices of $\epsilon_1$ and $B_*$
pay all these estimates. At entry $a_\psi=e^{\epsilon_1E_{\rm post}}$.
By the scaled-axis estimate of the radial lemmas, an extra
Euclidean Hessian with two further jets bounded by $D$ costs at
most $CD/(a_\psi H)$ in scaled axes. Consequently at entry and
through the later aligned intervals,
\begin{equation}\label{assembly:aligned-extra}
 |j^{\leq2}G_e|_{\rm scaled}\leq\delta/H.
\end{equation}
Here we leave a fixed fraction of the $\delta$ budget for future
wells. The ratio $D_{\rm out}/a_\psi$ is exponentially small at
entry and only decreases later, since $a_\psi$ increases.
During an actual profile switch its new offset is accounted for
by Proposition~\ref{center:change}; outside the switch, that offset
is constant.

\begin{lemma}\label{assembly:change-order}
At constant $E_{\rm post}$ the center can be changed from $\zeta_j$
to $\zeta_{j+1}$, and the data can then be restored to
\[
 \lambda=p=1,\qquad \kappa=1,\qquad H=\sigma,
\]
using finite clock intervals and only the aligned and center-change
regimes. The old-twist bounds needed for choosing the profile
shift are available independently of the targeted $\lambda,H$.
\end{lemma}

\begin{proof}
Initially grow $H$ with $\kappa=\kappa_*$ and $p=\lambda=1$.
Use the $\lambda=1$, $H\geq\sigma$ alternative of
Proposition~\ref{radial:aligned}. Ultimately we will lower $p$
slightly below one for a bump, thereby increasing $\lambda$, and
then restore $p=1$ and return smoothly to $\kappa=1$.

We first justify the order in which the targets are chosen.
During the growth and the bump, $p\leq1$ and $\kappa=\kappa_*$,
so $d\log H/du_s\geq.08$. Thus
\[
 \int_{u_{\rm entry}}^\infty \frac{du_s}{H}
 \leq \frac{1}{.08H_{\rm entry}}.
\]
This bounds the total increase of $\Psi$ before the final transition
independently of the size of either target. The transition from
$\kappa_*$ to one takes a fixed $\log Y$ interval and leaves
$H$ nondecreasing, adding another fixed bound to $\Psi$.
Moreover $\kappa\geq1$ gives
$d\log Y/du_s\geq1$, hence
\[
 0\leq\nu-\nu_{\rm entry}
 =\int \frac{du_s}{Y}\leq Y_{\rm entry}^{-1}
\]
through the growth and holding intervals. The old nonradial
twists therefore have bounded values on one fixed compact
spatial interval, independently of the targets. Their scaled
derivatives are bounded by the fixed Euclidean bounds, since
the scaled coordinate radii are at most one. Derivatives of
$\Psi$ in these axes are $O(1/H)\leq O(1/\sigma)$.
The forward switch and center move each occupy a fixed amount
of the variable $\Psi$; these amounts can be included in the
same bounds. Include the reserved future value and jet allowance
as well. These are exactly old-twist bounds independent of
$\lambda,H$, as required by Proposition~\ref{center:change}.

Choose its shift $L_0$ from these bounds and $E_{\rm post}$.
This fixes its profile constants $C_f$ and a sufficient target
$\lambda$, and that target determines a sufficient target $H$.
This is the order of parameter selection in
Proposition~\ref{center:change}. The radial path reaches the targets
in the other order: grow $H$ first until it exceeds every required
multiple of
\[
 C_{E_{\rm post}}(1+\log\lambda)^2
 \quad\hbox{and}\quad C_f(1+\log\lambda)^2.
\]
Choose a smooth bump of $1-p\geq0$ with small amplitude and
small derivatives in $\log Y$, and with integral in $u_s$
equal to the desired $\log\lambda$. Such a bump exists by
making its middle plateau arbitrarily long; the endpoint ramps
can have fixed small amplitude and sufficiently long fixed
length, while the plateau length supplies the remaining integral.
The conversion $d\log Y/du_s=\kappa(1+Y^{-1})$ and its bounded
derivatives show that bounds through order two in $\log Y$
follow from the same bounds in $u_s$.
Throughout this bump $H$ increases, and the already chosen lower
bound for $H$ pays the entire range $0\leq\log\lambda(s)\leq
\log\lambda$. Apply the large-$H$ alternative of
Proposition~\ref{radial:aligned}.

Restore $p=1$ and lower $\kappa$ smoothly to one. The downward
transition can be made slowly enough in $\log Y$ that
$Y\kappa_Y\geq-\kappa$, with absolute derivative bounds.
Since $\kappa\geq1$ here, $H$ stays large. Now $\kappa=p=1$,
so both $H$ and $\lambda$ are constant. Perform the forward
profile switch, the center move, and the reverse switch of
Proposition~\ref{center:change}. Interpolate the two center
values by a flat-ended convex combination, keeping the center
in the unit disk. The old bounds and targets have been chosen
in the order just proved. The changes take a bounded interval
of $\Psi$, hence a finite interval of $u_s$ because
$d\Psi/du_s=1/H$ and $H$ is fixed, even though this interval
can be very long. Restore $q_0$ and retain the new constant
offset. The center is now $\zeta_{j+1}$. The offset may be arbitrarily
large, but it is now constant and changes no derivative of the twist. Its
value becomes part of the old data when the next testing height is
chosen, before the next accuracy parameter $c$.

Raise $\kappa$ to $\kappa_*$ again and use a bump with
$1<p<\kappa_*$ to bring $\log\lambda$ down to zero.
A small positive plateau and fixed ramps give any prescribed
finite integral of $p-1$; stop precisely when that integral
equals the existing $\log\lambda$. Thus $\lambda$ never falls
below one. Also
\[
 \frac{d\log H}{du_s}
 =\kappa_*-p+\frac{\kappa_*-1}{Y}>0,
\]
so the large-$H$ aligned condition remains valid. Restore $p=1$.

Finally reduce $H$ to $\sigma$. Transition $\kappa$ to $1/2$,
hold there for an adjustable interval, then return to one from
below. For $p=1$ the exact relation is
\begin{equation}\label{assembly:H-descent}
 \frac{d\log H}{d\log Y}=1-\kappa^{-1}.
\end{equation}
Its value is $-1$ on the middle plateau. The endpoint transitions
have fixed finite effects on $\log H$; the initial $H$ can be
made as large as needed. Choose the plateau length so that the
final value is exactly $\sigma$. Throughout this descent
$H\geq\sigma$, $\lambda=1$, and $1/2\leq\kappa\leq\kappa_*$,
so the $\lambda=1$ aligned alternative applies. Choose the
downward transitions slowly enough to preserve the same
derivative inequality as above.

All these intervals have finite $u_s$ or finite $\log Y$ length.
Equations~\eqref{assembly:clocks} show that $Y$ remains finite
on each finite $u_s$ interval and that the radius advances by
a finite positive amount. A putative pole of an indefinitely
continued accelerated equation is therefore never reached
during an episode. This proves the claim.
\end{proof}

\begin{lemma}\label{assembly:reset}
From the endpoint in Lemma~\ref{assembly:change-order}, severity
can be increased and a new cap started so that the next trigger
has \eqref{assembly:trigger-initial}, while the metric continues
to satisfy the quiet criterion. The cap can be extended to
arbitrarily large spatial radii.
\end{lemma}

\begin{proof}
The spatial scale $u_s$ may now be enormous, and $E_{\rm post}$ need
not dominate it. Proposition~\ref{quiet:profile} permits any fixed
holomorphic affine frame at the evaluation point. We first use
$V$-unit affine axes, where the curvature scale is the fixed $\sigma$
and the relevant jets are bounded independently of $u_s$. This will
justify the initial severity increase before arranging the size
inequality needed for the next trigger.

First hold $\kappa=p=\lambda=1$ and $H=\sigma$.
Let $E$ increase smoothly as a function of $u_s$, starting at
$E_{\rm post}$, with a fixed flat-ended transition to slope
$4C_\uparrow$. Its derivatives through order four in $u_s$
are bounded by constants depending on the fixed parameters,
not on $\sigma$ or the possibly enormous $u_s$ at entry.
Continue long enough that
\begin{equation}\label{assembly:catchup}
 E\geq2C_\uparrow u_s
\end{equation}
and any desired fixed threshold. The greater slope ensures
that this occurs after a finite interval.

We verify the quiet criterion even before \eqref{assembly:catchup}.
Use $V$-unit affine axes. The scaled-jet radial estimates give
metric jets of orders $r=1,2$ bounded by $CH^{r/2}$ in these
axes, and scaled derivatives of $u_s$ through order four are
bounded by an absolute constant. Therefore
\[
 |\partial^{[r]}E|_{V}\leq C H^{r/2},\qquad 1\leq r\leq4.
\]
The extras in \eqref{assembly:aligned-extra} satisfy the same
bounds after this normalization, and $G\geq V$.
For $\sigma$ small, the severity terms in $\mathcal N$ are
small relative to $G$, so
$h\asymp(1+\tau)V$, with first derivatives bounded in these
axes by $C(1+\tau)\sqrt H$.
Its connection is consequently $O(\sqrt H)$ uniformly in
$\tau$. The first, complex second, and pure covariant second
severity derivatives are small for every $\tau$.

The raw curvature margin on this hold is
$C^V(T,T)\geq c\sigma|T|^2$ in $V$-unit axes, because
$\kappa=1$ and $H=\sigma$. All other raw data are bounded there
by fixed constants. The previously chosen $E_{\rm post}$
exceeds the fixed profile and $\sigma$ thresholds, so it
pays this margin even when it is much smaller than the current
$u_s$. Thus there is no gap before severity catches up.

After \eqref{assembly:catchup}, keep the same slope while
turning $\kappa$ into a cap. Here is an explicit admissible
transition. If $Y_0$ is its starting clock value, set
\[
 A(Y)=1+Y_0+\int_{Y_0}^Y \eta_0(v/Y_0)\,dv,\qquad
 \kappa(Y)=\frac{A(Y)}{1+Y},
\]
where $\eta_0$ is smooth, equals one near $1$, decreases
nonnegatively to zero, and is zero for arguments at least two.
For $Y\leq Y_0$ use $A=1+Y$.
Then $1\leq A\leq1+Y$, so
$(1+Y)^{-1}\leq\kappa\leq1$. Moreover
\[
 Y\kappa_Y+\kappa
 =\frac{YA'(Y)}{1+Y}+\frac{A(Y)}{(1+Y)^2}\geq0.
\]
The first two derivatives of $\kappa$ in $\log Y$ have absolute
bounds, because $Y/Y_0$ ranges over a fixed interval while
$\eta_0$ changes. For example, writing $b=Y/(1+Y)$ and
$D=Y\partial_Y$, one has
$D\kappa=b(\eta_0-\kappa)$ and
$D^2\kappa=b(1-2b)(\eta_0-\kappa)+bD\eta_0$.
At the end $A$ is the new constant $C_\kappa$.
Since $\kappa\leq1$ and $p=1$, $H$ is nonincreasing and
stays at most $\sigma$.

The same normalized jet and connection estimates continue
through the transition and the cap. To check the raw curvature
margin there, note that $Y\geq t$: it holds initially, and
$d\log Y/d\nu=\kappa(1+Y)\geq1$ preserves it.
Also $u_s\geq t$ after choosing the initial constant large
enough, since $du_s/dt=Y/t\geq1$.
Hence
\[
 \kappa H\geq c/(ta)\geq c e^{-2u_s}
 \quad(t\geq1),
\]
and \eqref{assembly:catchup} with large $C_\uparrow$ pays the
quiet lower bound $e^{-\epsilon E}$. The positive-order
extras stay bounded relative to $V$ on an arbitrarily long
cap. More explicitly, a fixed Euclidean jet bound $D_{\rm out}$
costs at most
$C(D_{\rm out}/a)H^{r/2}$ in $V$-unit axes for Hessian jets
of order $r\leq2$, since $aH=Y/t\geq1$.
The ratio only decreases as $a$ grows.
Future wells are included in the fixed allowance until
their own activations.

Coast as far in $\nu$ as desired. Shortly before the next
chosen trigger, turn off the severity slope smoothly over
a fixed-length $u_s$ interval. The margin in
\eqref{assembly:catchup} is sufficient to retain
$E_0\geq C_\uparrow u_s(\nu_B)$ after this bounded final
increment, by making the cap endpoint large enough.
The existing raw jet bounds can also be absorbed by making
that endpoint large. This produces
\eqref{assembly:trigger-initial} and the next cap data.
\end{proof}

\subsection{Global realization and completeness}

We now carry out the recursion. At each stage first extend the
cap so that its trigger starts beyond all previous stages and
with $\nu_B\geq j$. Fix its old-data bounds and testing height
by Lemma~\ref{assembly:height}. Then choose $c$ large enough
for \eqref{assembly:height-rate}, the polynomial approximation
for $h_j$, \eqref{assembly:diagonal}, and all the finite
thresholds in the trigger, test, post ramp, and aligned entry.
The later finite center-change targets are chosen in the
order proved in Lemma~\ref{assembly:change-order}.
Lemma~\ref{assembly:reset} completes the cycle.
Thus no clock or already configured profile is determined
by an unknown future well.

For completeness, the future allowance is valid in every
kind of local regime. In a quiet interval, later wells have
positive-semidefinite Hessians and bounded absolute jets;
they preserve $G\geq V$, and the preceding connection proof
uses exactly these bounds. In an aligned or switch interval,
their scaled Hessian jets are bounded by
$C/(a_\psi H)$, which fits the reserved part of
$\delta/H$ at the large chosen $a_\psi$. In a center move,
their values change the radius estimates by at most a fixed
factor, and their scaled derivatives are within the old-data
allowance included before choosing the shift. On a shear
collar they lower the twist by at most one, already included
in the lower bound used for its coefficient. Consequently
the local lemmas apply to the final infinite twist for
every fiber radius, even though curvature margins need not
have a positive infimum over all fiber radii.

The condition $\nu_B\geq j$ rules out finite accumulation of
stages. Any compact spatial set lies before the activation
of all sufficiently late wells. For a fixed derivative order
$r$, those later wells satisfy
$|\partial^{[r]}W_j|\leq2^{-j}$ once $j+4\geq r$.
The Weierstrass criterion in every derivative order proves
compact-smooth convergence of $\sum_jW_j$.
The $Z$ terms and nonconstant offsets are locally finite.
Flat-ended transitions, finite positive clock evolutions,
and the prescribed initial data therefore give global smooth
spatial data. On a compact spatial set the profile is smooth
to $\tau=0$ with $q(0,s)=1$; its normalized $F-x$ is smooth
there as well. The smooth inversion at the fiber center
therefore gives a global smooth potential
\[
 \ell(s,y)=\ell_{\rm rad}(s,y)+
 \operatorname{Re}\bigl(c_{\rm sh}(s)(y-\zeta(s))\bigr).
\]
The shears are absent where the center changes.
All local curvature and positivity conclusions apply,
including the fiber axes by their smooth limits.

We spell out the metric comparison used for completeness.
In the constant-center regimes the Schur metric is bounded
below by a fixed positive multiple of $V$; during profile
switches, positivity of $\mathcal N$ for all $\tau$ also
gives $h=V+\int_0^\tau\mathcal N\,dv\geq V$.
The shear and moving-center comparisons can each be chosen
to lose at most a factor $1/2$. Their spatial supports are
separated, so these losses do not multiply over infinitely
many stages at a point. Hence a single $c_1>0$ satisfies
\begin{equation}\label{assembly:spatial-completeness}
 g_\ell(\dot s,\dot y)\geq c_1 V(\dot s,\dot s)
 \geq c_1a_{\rm init}|\dot s|^2
\end{equation}
everywhere.

A finite-length curve therefore has finite Euclidean
spatial variation and remains in a compact spatial set.
On such a compact set, for large $|y|$, the metric is also
bounded below by a fixed positive multiple of the frozen
log-block metric at the point. The tilt coefficients are
uniformly bounded in $\tau$ there. Indeed they are spatial
derivatives of $\psi-F$: on quiet intervals they use bounded
$F_E$; on switch intervals they use bounded $F_\beta$, which
is constant in $x$ past the fixed switch cutoff; and on a
center move the profile is spatially fixed apart from the
already bounded twist. There are only finitely many profile
episodes on the compact set, and the well sum is smooth.
Thus some compact-set constant $C_K$ bounds the tilt one-form.

Also $q\geq1$ and $F\leq x$, so at sufficiently large $|y|$
the fiber value satisfies $\tau\geq1$ and $P=\tau q\geq1$,
uniformly on the compact spatial set.
The frozen fiber differential is $dy/(y-\zeta(s))$; freezing
means that the center has its point value in this comparison.
The block bound and the triangle inequality imply along
the curve, when $|y|$ is large,
\[
 \frac{|dy|}{|y-\zeta(s)|}
 \leq C_K\,ds_{g_\ell}+C_K|ds|.
\]
Both terms on the right have finite integral by
\eqref{assembly:spatial-completeness}. Since $|\zeta(s)|<1$,
this bounds the total variation of $\log|y|$ outside a
fixed disk and prevents fiber escape. Thus no divergent
curve has finite length. Equivalently, a Cauchy sequence
lies in a compact Euclidean coordinate set, where the
smooth positive metric gives convergence. The base metric
is complete.

\begin{proof}[Proof of Proposition~\ref{assembly:base}]
The recursion gives a smooth strictly plurisubharmonic
potential on $\mathbb C^{n+1}$, with nonpositive real
sectional curvature by the local lemmas and completeness
by the preceding argument. Every required pair
$(\zeta,h)$ is tested infinitely often. The stages satisfy
$R\geq1$, $N\to\infty$, $c\to\infty$,
$\log N\leq.005c<c/100$, $D_j\to\infty$, and
$D_j/c_j\to0$. Lemma~\ref{assembly:height} supplies the
height bound. The node disks, sphere derivative bound,
and shear amplitude have the fixed exponents
\eqref{assembly:test-exponents}. Their phases were chosen
as permitted by Proposition~\ref{grid:criterion}, and
their geometric realization allows all those choices.
That proposition proves the asserted Liouville property.
\end{proof}

\begin{remark}\label{assembly:potential-unbounded}
The potential is necessarily unbounded below, and this can also be
seen directly from the disks. At a testing node choose
$y-\zeta_j=-\tfrac12 e^{m_0c}\overline{b_v}$.
The radial increment is at most one, whereas the shear contributes
$-\tfrac12\alpha e^{m_0c}$. Thus
\[
 \ell(s,y)\leq D_j+1-\tfrac12e^{(m_0-p_1)c}\longrightarrow-\infty.
\]
Here $m_0>p_1$ and $c/(D_j+1)\to\infty$.
In particular the elementary holomorphic functions $e^{ikw}$,
$k>0$, are unbounded on the tube; they do not contradict its
bounded-function conclusion.
\end{remark}

\begin{proof}[Completion of the proof of Theorem~\ref{thm:main}]
Apply Lemma~\ref{tube:transfer} to the base just
constructed. Its tube is connected, complete, and
diffeomorphic to Euclidean space, hence simply connected.
Its real sectional curvatures are at most $-1$ by that lemma,
and its bounded holomorphic functions are constant by
Proposition~\ref{assembly:base}. The complex dimension is the fixed finite
$m=n+2$.

There is no finite lower sectional-curvature bound.
At $s=0,y=\zeta_1$, the center and profile are constant
nearby and there is no shear. On the center axis the
spatial curvature tensor of the base equals that of $V$:
the added radial-fiber potential vanishes to second
order in $y-\zeta_1$, and its cubic derivatives with two
spatial indices and one fiber index vanish there.
The initial radial metric has a strictly negative
holomorphic-plane curvature in a spatial direction.
Corollary~\ref{tube:negative-plane} therefore shows that the tube's
sectional curvatures are unbounded below.
This proves every assertion of Theorem~\ref{thm:main}.
\end{proof}


\paragraph{The two-sided question.}
Neither this example nor the independent pinched threefold of
\citet[Theorem~1.1]{OpenAIPinchedThreefold2026} decides whether a
nonconstant bounded holomorphic function must exist under finite
two-sided negative real sectional pinching. The present curvature is
unbounded below, while the threefold retains nonconstant bounded
base-coordinate functions.
